\chapter{Охлаждение в решётке}\label{ch:ch3}

Третья глава посвящена разработке протокола подготовки атомов на основном, а также на двух возбуждённых колебательных подуровнях оптической решётки.
Разрабатываемый протокол должен решать несколько проблем, стоящих перед нами после перезахвата атомов из магнито-оптической в оптическую решётку.
Во-первых, охлаждение атомного ансамбля до основного колебательного состояния оптической решётки. 
Существует несколько методов решения этой задачи.
В связи с тем, что расстояние между колебательными уровнями ионов в радиочастотных ловушках и атомов в оптических решетках (вдоль оси решетки) составляет порядка $100\unit{кГц}$, наиболее широко применяемым методом глубокого лазерного охлаждения является метод охлаждения на колебательных частотах (sideband cooling) \cite{neuhauser1978optical, perrin1998sideband}. 
В его основе лежит возбуждение атомного перехода с уменьшением колебательного квантового числа на единицу (красного сайдбенда) и последующий возврат атома в начальное состояние посредством спонтанного распада без изменения колебательного числа (вследствие работы в режиме Лэмба-Дике переходы с изменением колебательного числа оказываются менее  вероятными, чем без изменения). 
Можно выделить три схемы охлаждения на колебательных частотах: (1) с использованием перехода с естественной шириной $\sim 1\div10\unit{кГц}$, которая обеспечивает достаточно быстрый спонтанный распад \cite{ushijima2018operational};%накидать ещё
(2) рамановский переход с использованием различных уровней сверхтонкой структуры либо магнитных подуровней одного уровня \cite{chen2017sympathetic, seck2016raman}; 
(3) возбуждение узкого часового перехода на красном колебательном сайдбенде с последующей перекачкой населенности \cite{zhang2022subrecoil}. 
В схемах (2) и (3) для возвращения атома в начальное состояние используется перекачивающее излучение, возбуждающее атом на короткоживущий уровень.
Хотя для атомов тулия разрабатывался метод считывания населённости часовых подуровней с помощью дополнительного перекачивающего излучения \cite{tregubov2020detection}, но в дальнейшем от него отказались, и сейчас перекачивающее излучение в эксперименте отсутствует. 

\begin{figure}[t!]
=	\includegraphics[width=0.8\linewidth]{Chapter0/ОписаниеОхладженияСайдбэндов.pdf}
	\caption{
			Различные схемы лазерного охлаждения в оптической решётке до основного колебательного состояния.
			а) Возбуждение красного колебательного сайдбэнда часового перехода с последующей перекачкой атомов в основное состояние через промежуточный короткоживущий уровень. 
			Перекачка необходима так как время жизни возбуждённого состояния слишком велико для эффективного охлаждения.
			б) Охлаждение на спектрально-узком переходе, чья спектральная ширина с одной стороны позволяет селективное возбуждение красного колебательного сайдбэнда (красная линия), и с другой стороны, обеспечивает быстрый спонтанный распад без изменения колебательного числа (черная линия).
			в) Рамановский переход между двумя (магнитными) подуровнями основного состояния с уменьшением колебательного числа на единицу.
			г) Характерный спектр перехода в режиме разрешённых боковых полос для двух разных температур атомного ансамбля.
			Пунктирной линией показан спектр для случая $T_z = T_r = 10\unit{мкК}$ - характерной температуры атомного ансамбля после перезахвата атомов из МОЛ в оптическую решётку.
			Число атомов в основном колебательном состоянии определяет относительные интенсивности синего и красного сайдбэндов.
			Сплошная линия соответствует охлаждённому ансамблю атомов с $T_z = T_r = 1\unit{мкК}$.
			Большая часть атомов в этом случае расположена в основном колебательном состоянии, что приводит к почти полному исчезновению красной боковой полосы. 
		} 
	\label{fig:МетодыОхлаждения}
\end{figure}

В случае использования перехода на длине волны \thd, циклической компонентой является переход между крайними магнитными подуровнями двух состояний: $\ket{g, m_F = -4} \xrightarrow{} \ket{e, m_F = -5}$. 
Это приводит нас к необходимости предварительно осуществить накачку атомов на магнитный подуровень $\ket{g, m_F = -4}$. 
После охлаждения до основного колебательного состояния атомы остаются сосредоточенными на крайнем магнитном подуровне, а значит необходимо также решать проблему возврата части атомов на центральный магнитный подуровень и вычищения поляризации атомного ансамбля.
Наконец, либо до, либо после перекачки атомов в начальное состояние часового перехода требуется осуществить перенос населённости с основного на первое либо второе возбуждённое состояние. 
В данной главе будут решены все эти задачи. 

\section{Оптическая накачка}

После перезахвата атомов из магнито-оптической ловушки в оптическую решётку распределение населённостей магнитных подуровней хотя и смещено в сторону отрицательных проекций полного момента, однако далеко от чисто поляризованных атомов (см. рис \ref{fig:Накачка}), поэтому перед реализацией охлаждения до основного колебательного состояния необходимо сконцентрировать атомы на крайнем магнитном подуровне. 
Существует несколько способов решения этой задачи, разработанных как для атомов тулия, так и для других элементов. 
Как уже упоминалось ранее в главе \ref{ch:ch1}, при охлаждении атомов в магнито-оптической ловушке с большой отстройкой от резонанса, ансамбль атомов окажется в основном на крайнем магнитном подуровне.  
Такой метод используют, наши коллеги из РКЦ \cite{tsyganok2018polarized}. 
Как, опять же, упоминалось ранее, в таком случае геометрический объём МОЛ окажется слишком большим для эффективного перезахвата атомов в оптическую решётку, поэтому нам такой метод не подходит.
Оптическая накачка циркулярно-поляризованным излучением в тёмное состояние (например используя накачивающий переход на длине волны $418\unit{нм}$) также может решать эту задачу, но для её реализации требуется дополнительной лазерной системы, что усложняет установку, и чего хотелось бы избежать.
%Накачка последовательностью радиочастотных импульсов осложняется тем, что такие переходы очень уж чувствительны к магнитному полю. 
%Другие методы мне пока не известны, но я их обязательно заботаю и раскритикую.
Накачку циркулярно-поляризованным излучением можно провести и без наличия тёмного состояния, используя тот факт, что возбуждение циклического перехода c конечного состояния не меняет распределение атомов по колебательным подуровням.
Для реализации последнего метода можно использовать переход на длине волны \thd, который обладает подходящей для этого метода структурой уровней.

В прошлой главе было показано, что существует два режима рабочих параметров для работы с переходом на длине волны \thd в оптической решётке. 
Для подготовки атомов в состоянии \gmF{4}{-4} был использован режим слабых магнитных полей (магнитное поле $B = 150\unit{мГс}$ и глубина оптической решётки $U = 1000\Erec$), в которым $\sigma^-$ переходы могут возбуждаться одним мощным импульсом (см. рисунок \ref{fig:Накачка} в)). 
Оптическая накачка атомов на крайний магнитный подуровень осуществлялась пучками МОЛ, общая мощность излучения в эксперименте составляла $P \approx 250\unit{мкВт}$ ($s\approx30$), при длительности накачивающего импульса $\tau=50\unit{мс}.$
Чистота итоговой поляризации атомов определяется в эксперименте Штерна-Герлаха и составляет $\eta = 98(2)\%$, что объясняется как нерезонансным возбуждением $\pi$-переходов, ненакаченными атомами на $\ket{g, F = 4, m_F = 3,4}$, так и эффектом перемешивания состояний (Состояние $\ket{e, F = 5, m_F=-5}$ имеет ненулевую примесь состояния $\ket{e, F = 5, m_F=-3}$, из-за чего может распадаться на уровни $\ket{g, F = 4, m_F = -3(-2)}$). 
%Посчитать равновесное распределение, МБ для параметров примерно так и получится.

В ходе дальнейшей работы было сделано два улучшения этого метода. 
Во-первых, для уменьшения влияния эффекта подмешивания состояний, можно добавить дополнительный накачивающий импульс при большем магнитном поле. 
Примесь состояний уменьшается при увеличении магнитного поля, а подобрав его величину так, чтобы резонансная частота $\sigma^-$ переход с состояния $\ket{g, m_F = -3}$ будет совпадать с охлаждающей частотой для $\sigma^-$ перехода с состояния $\ket{g, m_F = -4}$ можно также сократить время, требуемое для дальнейшего охлаждения.
Так, для глубины оптической решётки $U=1000\Erec$ магнитное поле, для которого выполняется это условие составляет $B = 430\unit{мГс}$ (хотя на этом магнитном поле для крайнего магнитного подуровня наблюдается резонанс Фешбаха, время, необходимое для вычищения поляризации таким образом составляет порядка $\tau \sim 10\unit{мс}$, за которое не наблюдается больших потерь атомов).
Во-вторых, для уменьшения нерезонансного возбуждения $\pi$-переходов, можно воспользоваться эффектом подмешивания состояний для оптической накачки.
Т.к. для оптической накачки используются пучки МОЛ в области захвата атомов присутствует излучение всех поляризаций, что приводит к возможности возбуждения, например, $\sigma^{2-}$ переходов. 
Для таких переходов состояние $\ket{g, F = 4, m_F = -4}$ является тёмным, а значит попадая в него, атомы перестают взаимодействовать с излучением, что позволяет существенно уменьшить долю деполяризованных из-за эффекта подмешивания состояний атомов.
Дополнительным плюсом этого метода является то, что вместо нерезонансного возбуждения $\pi$ переходов, приводящего к деполяризации атомов, будет происходить нерезонансное возбуждение $\sigma^-$ переходов, что только увеличит скорость накачки.   
%Нарисовать ли картинку для этих случаев тоже?

\begin{figure}[htbp!]
	\begin{center}
		\includegraphics[width=\linewidth]{Chapter3/3_1_Накачка.pdf}
	\end{center}
	\caption{
		(а) Фотография облака, полученная в эксперименте Штерна-Герлаха после перезахвата атомов из МОЛ в оптическую решётку. Аппроксимация изображения выполнена суммой гауссиан, соответствующих атомам на всех подуровнях основного состояния. Атомы практически равномерно распределены по подуровням $m_F = -4..0$, в то время как на подуровнях $m_F = +3, +4$ атомы практически отсутствуют.
		(б) Фотография облака, полученная в эксперименте Штерна-Герлаха после оптической накачки. 98\% атомов сосредоточены на крайнем магнитном подуровне.
		(в) Частоты переходов в режиме малых магнитных полей ($B = 150\unit{мГс}$ выбрано для реализации оптической накачки).
		Сплошными линиями обозначены частоты $\sigma^-$ переходов, пунктирными линиями -- частоты $\pi$ переходов.
		Цветом показаны магнитные подуровни основного состояния, с которых происходит переход.
		В таком режиме большая часть $\sigma-$ переходов может возбуждаться одним мощным импульсом.
		Например, изображённый на фоне профиль линии соответствует лазерному излучению, уширенному мощностью с параметром насыщения $s = 20$.
	}
	\label{fig:Накачка}
\end{figure}


\section{Охлаждение}

Этот пункт будет посвящён охлаждению атомов тулия, накаченных на крайний магнитный подуровень $\ket{g, m_F = -4}$ до основного колебательного состояния оптической решётки. 
В первом подпункте приведён метод измерения распределения атомов по колебательным подуровням с помощью спектроскопии боковых колебательных частот часового перехода.
Второй подпункт посвящён процессу охлаждения атомов и включает в себя обоснование выбора параметров эксперимента, сам процесс охлаждения и полученные результаты. 

\subsection{Измерение распределения атомов по колебательным подуровням}

Наиболее распространённым методом измерения температуры атомов в оптической решётке является спектроскопия часового перехода в режиме разрешённых боковых колебательных частот (сайдбендов) \cite{}.
Простейшая модель предствляет потенциал решётки \ref{eq:LatticePotential}
в виде ангармонического осциллятора \ref{eq:LatticePotentialDecomposition}, после чего выводятся все необходимые соотношения.
Для больших подробностей нужно будет проконсультироваться с Костей, однако прямо сейчас можно будет записать основные соотношения. 
Для аппроксимации спектра синего сайдбенда используется следующее выражение для вероятности возбуждения $P_e$ в зависимости от отстройки $\delta>0$ от несущей используется формула \ref{eq:BlueSidebandProbability}, в которой $A_{n_z}$ - искомые населённости колебательных подуровней.
Использование собственных значений энергии $U_{n_z (r)}$,  полученных в приближении Борна-Оппенгеймера \cite{beloy2020modeling} позволяет вычислять в правой части взвешенные больцмановскими факторами вклады от атомов с различными значениями радиальных энергий $E_r$ для радиальной температуры $T_r$.
Для них находятся резонансные отстройки синего сайдбенда  $\nu^{n_z, r}_{bsb} = (U_{n_z + 1}(r) - U_{n_z}(r))/h$ параметры Лэмба-Дике $\eta^2 (n_z, r) = \nu_{rec}/\nu^{n_z, r}_{bsb}$ и частоты Раби $\Omega_{n_z, r}$ переходов $n_z \xrightarrow{} n_z + 1$ \ref{eq:BlueSidebandRabiFrequency}.
В этой формуле $L^m_n (x)$ - обобщённый полином Лагерра. 
Для красного сайдбенда верны те же формулы, с вычислением $\Omega_{n_z - 1, r}$ и $\nu^{n_z, r}_{rsb} = \nu^{n_z - 1, r}_{bsb}$, отрицательной отстройкой и удалением слагаемого с $n_z = 0$. 
Применимость объясняется тем, что величина параметра насыщения большая, как и длительность импульса.
Результатом аппроксимации спектра являются радиальная температура $T_r$, глубина решётки $u_lat$ положения резких склонов сайдбендов. 
В случае странных распределений атомов по колебательным подуровням следует предварительно найти "честную" глубину решётки. 
Модельные спектры боковых частот часового перехода для ансамблей атомов с различным распределением по колебательным подуровням представлены на рисунке \ref{fig:МодельныеСайдбенды}. 

\begin{equation}
	\label{eq:LatticePotential}
	U(z,r) = -U_0 \times \cos^2(kz)e^{-2r^2/\omega_0^2} + mgz
\end{equation}

\begin{equation}
	\label{eq:LatticePotentialDecomposition}
	U(z,r) = U_0(-1 + k^2 z^2 + ...) 
\end{equation}

\begin{equation}
	\label{eq:BlueSidebandProbability}
	P_e(\delta) = \frac{1}{2} \sum_{n_z} A_{n_z} \sum_{r} \frac{W_{n_z}(E_r, T_r)\Omega_{n_z, r}^2}{\Omega_{n_z, r}^2 + ...} 
\end{equation}

\begin{equation}
	\label{eq:BlueSidebandRabiFrequency}
	\Omega_{bsb} (n_z, r) = \Omega_0 e^{\eta^2 (n_z, r)/2} \eta (n_z, r) \sqrt{\frac{n_z !}{(n_z + 1)!}} \times L^1_{n_z} (\eta^2 (n_z, r))
\end{equation}

\begin{figure}[htbp!]
	\begin{center}
		\includegraphics[width=0.8\linewidth]{Chapter3/МодельныеСайдбенды}
	\end{center}
	\caption{
		Модельные спектры боковых частот полученные по формулам выше.
		Сплошной линией показаны атомы охлажденные до основного колебательного состояния.
		Пунктиром показан спектр атомов, находящихся в первом возбуждённом колебательном состоянии.
		Для сравнения, приведёно спектр боковых частот для ансамбля атомов, со средним колебательным числом $n_z = 1$ и радиальной температурой $T_r = 7\unit{мкК}$
	}
	\label{fig:МодельныеСайдбенды}
\end{figure}

\subsection{Охлаждение до основного колебательного состояния}

Оптическая накачка атомов на крайний магнитный подуровень основного состояния проводилась при магнитном поле $B = 150\unit{мГс}$, однако опрос часового перехода должен происходить при большем значении магнитного поля для уменьшения шумов, связанных с его флуктуациями, а также для уменьшения масштаба эффекта подмешивания состояний (см рис.\ref{fig:StateMixingSynthetic}).
При изменении магнитного поля может происходить нагрев атомов при проходе через резонанс Фешбаха, а значит охлаждение должно осуществляться на том же магнитном поле, на котором будет опрашиваться часовой переход.
Так как существуют резонансы Фешбаха для крайнего магнитного подуровня при $B \approx 430\unit{мГс}$ и $B \approx 800\unit{мГс}$ \cite{khlebnikov2019random}, было принято решение использовать лежащее между ними магнитное поле $B = 600\unit{мГс}$.
Глубина решётки в эксперименте поддерживалась постоянной и равнялась $U = 1000~\Erec$.

Для реализации охлаждения на боковых частотах перехода $\ket{g, m_F = -4} \xrightarrow{} \ket{e, m_F = -5}$ требуется построить модель, предсказывающую частоты переходов между различными колебательными подуровнями основного и возбуждённого состояния в зависимости от глубины решётки. 
Что-то похожее было сделано в предыдущем подпункте для часового перехода, однако для охлаждающего перехода длина волны оптической решётки хотя и близка к магической, но таковой всё таки не является. 
Приходится использовать поляризуемости, измеренные ранее для построения такой модели.
Результаты этого представлены на рисунке \ref{fig:CoolingSidebandModel}.
Можно выделить несколько характерных особенностей этой модели.
Во-первых, видно, что всё ещё присутствует характерный спектр боковых колебательных частот, что позволяет осуществлять эффективное охлаждение в аксиальном направлении оптической решётки одним мощным охлаждающим импульсом. 
В радиальном же направлении работает обычное доплеровское охлаждение, т.к. переходы без изменения колебательного числа (на несущей) расположенны компактно, близко друг к другу, что позволяет охлаждать в радиальном направлении атомы, располагающиеся на всех колебательных подуровнях одним  охлаждающим импульсом, частота которого отстроена в красную сторону от резонанса.
Во-вторых, при используемых в эксперименте глубинах оптической решётки число уровней энергии составляет $N \approx 15$, при этом, частоты Раби переходов с уменьшением колебательного числа не обращаются в ноль, как бывает в экспериментах по охлаждению ионов в радиочастотных ловушках.
В-третьих, можно использовать ``вычищение атомов'' по методике описанной ниже. 

\begin{figure}[htbp!]
	\begin{center}
		\includegraphics[width=0.8\linewidth]{Chapter3/4СпектрСайдбэндов506}
	\end{center}
	\caption{
		Показаны частоты переходов между различными колебательными подуровнями для двух разных компонент перехода (охлаждающий $\ket{g, m_F = -4} \xrightarrow{} \ket{e, m_F = -5}$, и переход, который мог бы использоваться для удаления горячих атомов c $m_F = 0$: $\ket{g, m_F = 0} \rightarrow{} \ket{e, m_F = -1}$).
		Частоты приведены для условий эксперимента: глубина решётки $U = 1000\text{ E}_{rec}$, магнитное поле $B=600\unit{мГс}$.
		Частоты на графиках отсчитываются от частоты невозмущённого перехода.
		Красным цветом обозначены частоты переходов с уменьшением колебательного числа на единицу, зелёным -- без изменения колебательного числа, синим -- с увеличением на единицу.
		Оттенок цвета показывает колебательный подуровень с которого возбуждается переход.
		(а) Для охлаждающего перехода показано доплеровское уширение до охлаждения (светло-серый контур), и после охлаждения (тёмно-серый контур), соответствующее температуре $T_r = 13\unit{мкК}$ и $T_r = 3\unit{мкК}$ соответственно.
		Чёрным пунктиром показаны частоты импульсов охлаждающего излучения. 
		(б) Частоты переходов между различными колебательными подуровнями оптической решётки для перехода с центрального магнитного подуровня. 
		Эффективное охлаждение аналогичное (а) здесь было бы невозможно. 
		Однако, возможно финальное вычищение атомов, находящихся не на нулевом колебательном подуровне: возбуждая переходы с уменьшением колебательного числа мы добьёмся того, что все атомы не на нулевом колебательном подуровне в потенциале оптической решётки либо попадут на основное колебательное состояние, либо покинут центральный магнитный подуровень и будут удалены из решётки в процессе очищения поляризации атомов.
	}
	\label{fig:CoolingSidebandModel}
\end{figure}

Для реализации охлаждения предполагалось использовать пучки МОЛ, однако в данном эксперименте оказалась важной разница геометрий пучков в двух системах.
Так, в большой системе пучок, близкий к вертикальному, эффективно охлаждал атомы вдоль оси решётки, в то время как два других позволяли охлаждать атомы в радиальном направлении.  
В компактной же потребовался дополнительный охлаждающий пучок вдоль оси решётки, т.к. при взаимодействии с излучением под углом $45^\circ$ к вертикали, при поглощении фотона из него происходил нагрев атома в одном направлении при охлаждении в другом.

Охлаждение осуществляется двумя импульсами длительностью $\tau = 50$ мс для охлаждения в радиальном и аксиальном направлениях, соответственно.
Характерные отстройки частоты от невозмущенного резонанса для каждого из импульсов приведены на рис. \ref{fig:CoolingSidebandModel} и составляют: для охлаждения в радиальном направлении $\approx -1560\unit{кГц}$, для охлаждения в аксиальном направлении $\approx -1590\unit{кГц}$, что соответствует отстройкам $\approx -40\unit{кГц}$ и $\approx -70\unit{кГц}$ от частоты несущей.
Характерные мощности для этих импульсов составляют около $200\unit{мкВт}$, что соответствует параметру насыщения $s \approx 25$.
Полученный в результате такого охлаждения ансамбль атомов охлаждён до основного колебательного состояния в оптической решётке со средним колебательным числом $\bar n_z = 0.01(1)$ и радиальной температурой $T_r \approx 3 \unit{мкК}$.

Спектры боковых частот для измерения распределения атомов по магнитным подуровням измерялся для перехода $\ket{g, m_F = -4} \xrightarrow{} \ket{c, m_F = -3}$. 
Т.к. поляризуемости ``часового'' состояния известны, и полная поляризуемость верхнего уровня не слишком сильно отличается от нижнего, модель спектра боковых колебательных частот часового перехода нуждалась лишь в небольшой модификации. 
Результаты представлены на рисунке \ref{fig:Охлаждение} (б) - г)). 
Достигнутое в результате среднее колебательное число $\bar n_z = 0.01(1)$, ограничивается, по всей видимости, нерезонансным возбуждением несущей частоты вместо красного сайдбенда. %посчитать не должно быть сложно
Также, это ограничение может быть связано с пределами считывания камеры.
Такое значение находится на уровне лучших достигнутых ансамблей в оптической решётке \cite{}.
Отметим, что если осуществить радиальное охлаждение во время оптической накачки, и использовать дополнительные импульсы для вычищения поляризации, как было указанно в конце прошлого пункта, то пропадает необходимость в радиальном охлаждении атомов, а время, необходимое для охлаждения в аксиальном направлении сокращается до $\tau = 5\unit{мс}$. 

\begin{figure}[htbp!]
	\begin{center}
		\includegraphics[width=0.8\linewidth]{Chapter3/3_2_НастоящиеСайдбэнды.pdf}
	\end{center}
	\caption{
		Результаты измерения температуры атомов в оптической решётке методом спектроскопии боковых колебательных частот.
		а) Атомы перекаченные на центральный магнитный подуровень.
		Спектр соответствует двухтемпературному распределению: большая часть атомов (>95\%) находится в основном колебательном состоянии (соответствует ансамблю из пункта г)).
		Итоговое среднее колебательное число $\bar n_z = 0.07(2)$
		б) До охлаждения. 
		Спектр соответствует $\bar n_z \approx 2$ и радиальной температуре $T_r \approx 20\unit{мкК}$.
		в) Охлаждение в радиальном направлении.
		Спектр соответствует $\bar n_z = 0.76(6)$ и радиальной температуре $T_r = 6.4(2)\unit{мкК}$.
		г) Полное охлаждение.
		Спектр соответствует $\bar n_z = 0.01(1) $ и радиальной температуре $T_r = 4.1(1)\unit{мкК}$.
	}
	\label{fig:Охлаждение}
\end{figure}

\section{Перекачка атомов}

После реализации охлаждения в прошлом пункте, ансамбль атомов находится на крайнем магнитном подуровне, а значит, что для опроса часового перехода требуется перенос части населённости в состояние $\ket{g, m_F = 0}$.
В главе \ref{ch:ch0} описывался метод удаления атомов вне  центрального магнитного подуровня с помощью радиочастотного перехода между центральными магнитными подуровнями сверхтонких компонент основного состояния \cite{mishin2024combined}, поэтому достаточно будет перенести лишь часть атомов в это состояние. 
Такая задача не типична для оптических часов: в случае фермионных изотопов центрального магнитного подуровня не существует, и для компенсации эффекта Зеемана опрашиваются попеременно два крайних\cite{}, подготовить атомы на которых проще, а в случае бозонов центральный магнитный подуровень часто также является единственным \cite{}.

Перенос населённости на центральный магнитный подуровень можно осуществить попеременно возбуждая микроволновые переходы между различными магнитными подуровнями сверхтонкой структуры (например, подобное осуществлялось в работе \cite{}), однако сложность состоит в том, что матричные элементы у различных переходов могут сильно отличаться, и оказывается технически сложно возбуждать их все по очереди. 
В нашем случае это также затрудненно из-за одного на две системы генератора микроволнового сигнала и сложностей с выставлениями одинакового магнитного поля в двух системах и его контроле (это необходимо для равенства частот переходов в двух системах).
Перенос атомов также можно осуществить с помощью пары двухфотонных (Рамановских) переходов \cite{}, однако во-первых, снова встаёт вопрос контроля магнитного поля, а во-вторых, требуется снять вырождение между разными переходами, что в случае атомов тулия получается плохо из-за малой тензорной поляризуемости основного состояния. %описать подробнее с цифорками
Можно было бы использовать эффект подмешивания состояний для возбуждения, например $\sigma^{3+}$ переходов, однако при этом происходит нагрев атомного ансамбля. %попытаться показать, что происходит для короткого импульса при рассеянии одного фотона
В итоге было принято решение реализовать перенос атомов на центральный магнитный подуровень с помощью возбуждения радиочастотных переходов между зеемановскими компонентами основного состояния \cite{}.

Для задающего магнитного поля $B=600\unit{мГс}$ разница по энергии между двумя соседними магнитными подуровнями составляет $\Delta E = h\times 800\unit{кГц}$. 
При этом, оптическая решётка на длине волны $1063.5\unit{нм}$ снимает вырождение засчёт тензорной поляризуемости основного состояния, причём разница по частоте между переходами составляет до $30\unit{кГц}$.
Возбуждать $\pi$-импульсы тяжело из-за нерезонансных возбуждений, малой разницы по частоте между переходами, а также шумов магнитного поля, поэтому перекачка осуществлялась одним из двух вариантов.
В случае компактной системы, частота сканировалась через резонансы переходов с уменьшением амплитуда ближе к концу, а в случае большой системы, глубина оптической решётки уменьшалась так, что переходы становились вырожденными, и возбуждались все сразу. 
В итоге, для первого случая, удалось достичь населённости центрального магнитного подуровня $\eta\approx30\%$ от начального числа атомов, а во втором случае, удавалось достичь $\eta=1/9$ от начального числа атомов, при гораздо меньшей чувствительности этого числа к шумам магнитного поля.

Для вычищения поляризации используется возбуждение $\pi$-импульса микроволнового перехода между центральными магнитными подуровнями двух сверхтонких компонент основного состояния, после чего невозбуждённые атомы удаляются из оптической решётки импульсом охлаждающего излучения. 
Источником сигнала служит генератор SRS SG382, сам сигнал подаётся на самодельную антенну.
Оставшиеся на центральном магнитном подуровне состояния $\ket{g, F = 3}$ атомы могут быть далее перераспределены между уровнями $\ket{g, F = 3, m_F = 0}$ и $\ket{g, F = 4, m_F = 0}$ ещё одним импульсом микроволнового излучения.

Характерный спектр боковых частот для атомов, полученных таким образом приведён на рисунке \ref{fig:Охлаждение} а).

\section{Получение возбуждённых колебательных состояний}

После получения ансамбля атомов, охлаждённого до основного колебательного состояния, последней задачей этой главы является перенос населённости на возбуждённые колебательные подуровни. 
Существует несколько методов решения этой задачи. 
Так, в работе \cite{ushijima2018operational} предлагается использовать для этих целей быстрый адиабатический перенос населённости, при котором частота часового излучения сканируется через несущую и синий сайдбенд, а затем удаляются все невозбуждённые атомы. 
В нашем эксперименте тяжело было бы достичь таких интенсивностей насыщения, поэтому такой вариант нас не устраивает.
В иных случаях можно также как мы делали раньше калибровать ансамбль атомов после просеивания \cite{}, однако нас такое не устраивает, т.к. атомы уже на основном колебательном и у нас ничего не получится.  
Альтернативным вариантом является возбуждение $\pi$-импульса какого-либо перехода с изменением колебательного числа на единицу, что было продемонстрировано в работах \cite{}.
На этом варианте мы и остановились. 

Есть два перехода, которые позволяют осуществить это: часовой и охлаждающий.
Был проведён расчёт, показывающий, что для охлаждающего перехода получается либо возбуждать несущую нерезонансно, либо не выходит возбудить вдосталь атомов на возбуждённые колебательные. 
Поэтому был выбран вариант с часовым переходом (расчёты показали, что с ним всё ок).

Идея предлагаемого метода описана на рисунке \ref{} вкратце такова:  
\begin{itemize}
	\item Подготовка атомов на центральном магнитном подуровне в основном колебательном состоянии описанным ранее методом.
	\item Возбуждение $\pi$-импульса часового перехода с изменением колебательного числа на единицу.
	\item Удаление невозбуждённых атомов.
	\item Возврат населённости в основное состояние на несущей частоте часового перехода.
	\item Повторив это ещё один раз мы получим атомы на втором возбуждённом колебательном состоянии. 
\end{itemize}

Результаты применения этого метода показаны на рисунке \ref{fig:ExcitedMotionalSubstates}.
Для эксперимента по подготовке атомов на первом колебательном состоянии результаты следующие: максимально достижимая эффективность возбуждения атомов составляла $\eta=40\%$, при длительности импульса $\tau=330\unit{мкс}$. 
Интересный эффект состоит в уменьшении радиальной температуры после этого с $T_r=4.7\unit{мкК}$ до $T_r=2.1\unit{мкК}$, что объясняется разной эффективностью возбуждения для разных групп атомов.
Значение населённости, которое получается из снятого спектра составляли $p_1 = 95\%$, $p_0 = 4\%$. 
Для повторного применения этого метода результаты следующие: эффективность выше (в конце получаем $24\%$ от начального числа атомов), радиалка ещё ниже ($T_r = 1.8\unit{мкК}$), а чистота конечного состояния лучше (до 98\%).
%\textbf{П. 5} посвящён подготовке атомов на первых возбуждённых колебательных подуровнях основного состояния ($n_z = 1, 2$) 
%Для этого возбуждался $\pi$-импульс часового перехода с увеличением колебательного кванта на единицу, затем невозбудившиеся атомы удалялись из оптической решётки, а возбуждённые атомы возвращались в основное состояние $\pi$-импульсом на несущей частоте часового перехода.
%Таким образом, атомы оказывались на колебательном подуровне с $n_z= 1$.
%Повторение этого цикла ещё раз позволяет также получить атомы, находящиеся на колебательном подуровне $n_z= 2$.
%Населённость различных колебательных подуровней определялась по спектру часового перехода в режиме разрешённых боковых колебательных частот.
%Итоговая эффективность переноса атомов на первый и второй колебательные подуровни составляла $\eta = 40(24)\%$ соответственно, при доле атомов на целевом колебательном подуровне более 93\% (рисунок \ref{fig:ВозбуждённыеКолебательные}).

\begin{figure}[htbp!]
	\begin{center}
		\includegraphics[width=0.75\linewidth]{Chapter3/3_3_ВозбуждённыеКолебательные.png}
	\end{center}
	\caption{
		Спектры часового перехода в режиме разрешённых боковых колебательных частот для анасмблей атомов в различных колебательных состояниях оптической решётки.
		а) Атомы охлаждённые до основного колебательного состояния.
		б) Атомы, перенесённые на первое возбуждённое колебательное состояние ($n_z = 1$). 
		в) Атомы, перенесённые на второе возбуждённое колебательное состояние ($n_z = 2$). 
	}
	\label{fig:ExcitedMotionalSubstates}
\end{figure}


\section{Основные результаты главы \ref{ch:ch3}}

\begin{itemize}
	\item Получены ансамбли атомов тулия, захваченные в оптическую решётку с длиной волны $1063.5\unit{нм}$, со средним колебательным числом $n_z = 0.01(1)$.
	\item Получены ансамбли атомов тулия в возбуждённых колебательных состояниях оптической решётки с $n_z = 1$ и $n_z = 2$. 
	Населённость конечного колебательного состояния составляла более $90\%$.
\end{itemize}
 